\section{Contoh Terhitung}

\begin{example}
\textbf{Enkripsi dengan \textit{stream cipher}.} Plainteks dan keystream masing-masing 24 bit, ditulis dalam heksadesimal sebagai $P=\code{571964}$ dan $K=\code{F5923D}$. Nyatakan keduanya dalam biner, lalu XOR bit per bit, $c_i=p_i\oplus k_i$:
\[
\begin{array}{rcl}
P &=& 0101\,0111\,0001\,1001\,0110\,0100\\
K &=& 1111\,0101\,1001\,0010\,0011\,1101\\
\hline
C &=& 1010\,0010\,1000\,1011\,0101\,1001
\end{array}
\]
sehingga cipherteksnya adalah $C=\code{A28B59}$. Dekripsi memakai keystream yang sama:
\[
C\oplus K=\code{A28B59}\oplus\code{F5923D}=\code{571964}=P.
\]
Karena tiap bit diproses sendiri-sendiri, galat satu bit pada cipherteks hanya merusak satu bit plainteks hasil dekripsi: mengubah bit ke-5 pada $C$ menjadi $\code{AA8B59}$ memberi $P'=\code{5F1964}$, yang berbeda dari $P$ tepat pada bit ke-5 saja.
\end{example}

\begin{example}
\textbf{Membangkitkan keystream dengan LFSR.} Pakailah LFSR 4 bit pada
Bagian~\ref{sec:konsep-stream-cipher}, yaitu umpan balik $f=b_3\oplus b_0$
dengan muatan awal $1111$. Tabel~\ref{tab:lfsr-4bit} menunjukkan bahwa barisan
keluarannya berperiode 15 bit:
\[
\code{1111\,0101\,1001\,000}\ .
\]
Karena pesannya lebih panjang daripada periode itu, keystream diulang dari awal.
Untuk pesan \code{Aku} yang panjangnya 24 bit, keystream yang dipakai adalah
15 bit pertama ditambah 9 bit awalnya lagi:
\[
K=\code{1111\,0101\,1001\,0001\,1110\,1011}=\code{F591EB}.
\]
Enkripsi dilakukan dengan XOR per bit:
\[
\begin{array}{rcl}
P &=& 0100\,0001\,0110\,1011\,0111\,0101 \quad(\code{Aku})\\
K &=& 1111\,0101\,1001\,0001\,1110\,1011\\
\hline
C &=& 1011\,0100\,1111\,1010\,1001\,1110
\end{array}
\]
sehingga cipherteksnya adalah $\code{B4FA9E}$. Dekripsi memakai keystream yang
sama: $\code{B4FA9E}\oplus\code{F591EB}=\code{416B75}$, yaitu \code{Aku} kembali.

Perhatikan pengulangan yang tampak pada baris $K$: sembilan bit pertama
keystream muncul kembali pada sembilan bit terakhir, karena periode keystream
hanya 15 bit sedangkan pesannya 24 bit. Pengulangan inilah yang harus dihindari
oleh keystream generator yang baik, dan inilah alasan pembahasan pada
Bagian~\ref{sec:konsep-stream-cipher} menekankan pentingnya periode yang jauh
lebih panjang daripada pesan.
\end{example}

\begin{example}
\textbf{RC4 versi mainan ($N=8$) untuk kunci \code{AB}.} Agar seluruh langkah
dapat diperiksa dengan tangan, larik $S$ diperkecil dari 256 menjadi 8 elemen,
sehingga semua operasi modulo 256 menjadi modulo 8. Kunci \code{AB} memberi
byte $\code{41}_{16}=65$ dan $\code{42}_{16}=66$, yang setelah direduksi modulo
8 menjadi $K=[\,1,\ 2\,]$; kunci diulang periodik: $K_0=1$, $K_1=2$, $K_2=1$,
dan seterusnya.

\textbf{KSA.} Larik $S$ mula-mula berisi $[\,0,1,2,3,4,5,6,7\,]$. Penunjuk $j$
bergerak menurut $j \leftarrow (j + S_i + K_{i \bmod 2}) \bmod 8$, lalu $S_i$
dan $S_j$ dipertukarkan. Hasil setiap langkah adalah
\[
\begin{array}{cll}
i=0: & j = (0+0+1) \bmod 8 = 1 & \text{swap } S_0 \leftrightarrow S_1
      \rightarrow [\,1,0,2,3,4,5,6,7\,]\\
i=1: & j = (1+0+2) \bmod 8 = 3 & \text{swap } S_1 \leftrightarrow S_3
      \rightarrow [\,1,3,2,0,4,5,6,7\,]\\
i=2: & j = (3+2+1) \bmod 8 = 6 & \text{swap } S_2 \leftrightarrow S_6
      \rightarrow [\,1,3,6,0,4,5,2,7\,]\\
i=3: & j = (6+0+2) \bmod 8 = 0 & \text{swap } S_3 \leftrightarrow S_0
      \rightarrow [\,0,3,6,1,4,5,2,7\,]\\
i=4: & j = (0+4+1) \bmod 8 = 5 & \text{swap } S_4 \leftrightarrow S_5
      \rightarrow [\,0,3,6,1,5,4,2,7\,]\\
i=5: & j = (5+4+2) \bmod 8 = 3 & \text{swap } S_5 \leftrightarrow S_3
      \rightarrow [\,0,3,6,4,5,1,2,7\,]\\
i=6: & j = (3+2+1) \bmod 8 = 6 & \text{swap } S_6 \leftrightarrow S_6
      \rightarrow [\,0,3,6,4,5,1,2,7\,]\\
i=7: & j = (6+7+2) \bmod 8 = 7 & \text{swap } S_7 \leftrightarrow S_7
      \rightarrow [\,0,3,6,4,5,1,2,7\,]
\end{array}
\]
Dua langkah terakhir tidak mengubah apa pun karena indeks yang dipertukarkan
sama. Setelah KSA, larik $S$ menjadi $[\,0,3,6,4,5,1,2,7\,]$ --- inilah
permutasi yang dipakai PRGA.

\textbf{PRGA.} Penunjuk $i$ dan $j$ direset ke nol, lalu untuk setiap byte
keystream dihitung $i \leftarrow (i+1) \bmod 8$, $j \leftarrow (j+S_i) \bmod 8$,
pertukaran $S_i \leftrightarrow S_j$, dan $t \leftarrow (S_i+S_j) \bmod 8$.
Byte keystream adalah $S_t$. Untuk larik $S$ di atas hasilnya adalah
\[
\begin{array}{c|c|c|c|c}
\text{langkah} & i & j & \text{swap} & t \ \rightarrow\ u\\
\hline
1 & 1 & 3 & S_1 \leftrightarrow S_3 & (4+3) \bmod 8 = 7 \rightarrow S_7 = 7\\
2 & 2 & 1 & S_2 \leftrightarrow S_1 & (4+6) \bmod 8 = 2 \rightarrow S_2 = 4\\
3 & 3 & 5 & S_3 \leftrightarrow S_5 & (5+3) \bmod 8 = 0 \rightarrow S_0 = 0\\
4 & 4 & 7 & S_4 \leftrightarrow S_7 & (7+3) \bmod 8 = 2 \rightarrow S_2 = 4\\
5 & 5 & 0 & S_5 \leftrightarrow S_0 & (0+1) \bmod 8 = 1 \rightarrow S_1 = 6
\end{array}
\]
Jadi tiga byte keystream pertama adalah $7$, $4$, dan $0$. Bila plainteksnya
\code{Halo}, yaitu byte $\code{48}$, $\code{61}$, dan $\code{6C}$, maka
\[
\begin{array}{c|ccc}
\text{plainteks} & \code{48} & \code{61} & \code{6C} \quad(\text{\code{Halo}})\\
\text{keystream} & \code{07} & \code{04} & \code{00}\\
\hline
\text{cipherteks} & \code{4F} & \code{65} & \code{6C}
\end{array}
\]
Cipherteks \code{4F656C} tampak seperti data acak, dan dekripsi mengulang
prosedur PRGA yang sama persis untuk mengembalikan \code{Halo}.

Contoh ini sekaligus memperlihatkan mengapa ukuran larik $S$ yang sebenarnya
dipilih 256 dan bukan 8. Dengan $N=8$ hanya ada $8! = 40.320$ permutasi larik
yang mungkin, dan setelah KSA larik $S$ hasilnya dapat dicari dengan menelusuri
seluruh kemungkinan itu. Pada RC4 yang sesungguhnya, $N = 256$ sehingga jumlah
permutasi mencapai $256!$ --- bilangan dengan lebih dari 500 digit --- dan
penelusuran menyeluruh menjadi mustahil. Yang menjatuhkan RC4 bukanlah
kekurangan ukuran lariknya, melainkan korelasi antara byte awal keystream dan
kunci.
\end{example}

\begin{example}
\textbf{Satu \textit{quarter round} ChaCha20.} Kerjakan QR ChaCha20 atas
empat kata $a=1$, $b=2$, $c=3$, dan $d=4$ (semuanya dalam heksadesimal; tiap
kata adalah bilangan 32 bit, sehingga batas penjumlahannya adalah
$2^{32}=\code{FFFFFFFF}$). Urutan operasinya adalah penjumlahan, XOR, dan
rotasi seperti yang tertulis pada Bagian~\ref{sec:algoritma-stream-cipher}.
Delapan langkah yang harus dikerjakan adalah sebagai berikut:
\[
\begin{array}{lr}
a \leftarrow a+b                          & \code{00000003}\\
d \leftarrow \mathrm{rotl}(d \oplus a,16) & \code{00070000}\\
c \leftarrow c+d                          & \code{00070003}\\
b \leftarrow \mathrm{rotl}(b \oplus c,12) & \code{70001000}\\
a \leftarrow a+b                          & \code{70001003}\\
d \leftarrow \mathrm{rotl}(d \oplus a,8)  & \code{07100370}\\
c \leftarrow c+d                          & \code{07170373}\\
b \leftarrow \mathrm{rotl}(b \oplus c,7)  & \code{8B89B9BB}
\end{array}
\]
Sebagai pemeriksaan, kerjakan langkah ketiga dan keempat secara rinci. Pada
saat itu $a=\code{00000003}$ dan $d=\code{00070000}$, sehingga
\[
c = \code{00000003} + \code{00070000} = \code{00070003},
\]
lalu
\[
b = \mathrm{rotl}(\code{00000002} \oplus \code{00070003},\ 12)
  = \mathrm{rotl}(\code{00070001},\ 12)
  = \code{70001000}.
\]
Hasil akhir seluruh QR adalah $a=\code{70001003}$, $b=\code{8B89B9BB}$,
$c=\code{07170373}$, dan $d=\code{07100370}$.

Perhatikan dua hal. Pertama, masukan yang sangat sederhana ($1, 2, 3, 4$)
berubah menjadi empat kata yang tampak acak hanya dalam satu QR. Kedua,
$\mathrm{rotl}(x,n)$ berarti memutar bit $x$ ke kiri sejauh $n$ posisi: bit yang
keluar dari ujung kiri masuk kembali dari ujung kanan, berbeda dari pergeseran
biasa yang membuang bit tersebut. Pada rotasi 16 bit misalnya,
$\code{00000007}$ menjadi $\code{00070000}$ karena tiga bit terendah yang
bernilai $1$ berpindah tempat, bukan hilang.

Contoh ini menunjukkan mengapa ChaCha20 disebut memakai operasi AXR: seluruh
langkah di atas hanya memakai penjumlahan modulo $2^{32}$, XOR, dan rotasi ---
tidak ada perkalian, tidak ada tabel substitusi, dan tidak ada pembagian.
Konsekuensinya, satu QR dapat dikerjakan prosesor dengan beberapa instruksi
saja, dan itulah sumber kecepatan ChaCha20 pada perangkat lunak.
\end{example}
